%% guide-bac-centrage-tube — dispositif de centrage d'un tube sur une découpeuse laser,
%% vue de face en coupe schématique. Une seule vis (3), deux tronçons de filetage de sens
%% opposés, deux noix (4) et (6) guidées en translation par le bâti (1) : quand la vis
%% tourne, les deux mors en V (5) et (7) se rapprochent et centrent le tube T.
%% Chaque pièce a sa propre couleur : le regroupement en classes d'équivalence n'est pas
%% donné par la figure, il est à trouver.
%% \rappel cache l'encadré du bas : E0={1,2}, E1={3}, E2={4,5}, E3={6,7} (réponse de la Q1).
\documentclass[tikz,border=4pt]{standalone}
\usepackage{liaisons}
\begin{document}
\begin{tikzpicture}[x=1cm,y=1cm,
  piece/.style={line width=0.9pt, line join=round},
  fleche/.style={-{Stealth[length=6pt]}, line width=1pt},
  renvoi/.style={line width=0.4pt, draw=black!60},
  etiq/.style={font=\small, inner sep=2pt},
  num/.style={font=\small\bfseries, inner sep=1pt, fill=white}]

%% ---------------- 1 : le bâti (semelle rainurée + montant du palier droit) ----------------
\hachures[draw=black!75]{(0.30,-0.70) -- (0.30,0) -- (2.20,0) -- (2.20,-0.30)
  -- (7.20,-0.30) -- (7.20,0) -- (9.40,0) -- (9.40,-0.70) -- cycle}
\hachures[draw=black!75]{(8.95,0) rectangle (9.40,1.20)}
\node[num, text=black!75] at (1.05,-0.48) {1};

%% ---------------- 2 : le moteur ------------------------------------------------------------
\draw[piece, draw=solideE, fill=solideE!12] (0.55,0.05) rectangle (1.75,0.95);
\draw[piece, draw=solideE, line width=0.6pt] (0.72,0.05) -- (0.72,0.95) (0.89,0.05) -- (0.89,0.95);
\node[num, text=solideE] at (1.35,0.62) {2};

%% ---------------- 4 et 6 : les noix (pied engagé dans la rainure du bâti) -------------------
\draw[piece, draw=solideA, fill=solideA!10]
  (2.80,1.00) -- (2.80,0) -- (2.95,0) -- (2.95,-0.27) -- (3.65,-0.27) -- (3.65,0) -- (3.80,0) -- (3.80,1.00) -- cycle;
\node[num, text=solideA] at (3.30,0.78) {4};
\draw[piece, draw=solideF, fill=solideF!10]
  (5.60,1.00) -- (5.60,0) -- (5.75,0) -- (5.75,-0.27) -- (6.45,-0.27) -- (6.45,0) -- (6.60,0) -- (6.60,1.00) -- cycle;
\node[num, text=solideF] at (6.10,0.78) {6};

%% ---------------- 5 et 7 : les mors en V (contacts tangents au tube) -----------------------
\draw[piece, draw=solideC, fill=solideC!10]
  (2.45,1.00) -- (4.10,1.00) -- (4.10,1.35) -- (3.50,1.95) -- (4.10,2.55) -- (4.10,2.90) -- (2.45,2.90) -- cycle;
\node[num, text=solideC] at (2.90,1.95) {5};
\draw[piece, draw=solideG, fill=solideG!10]
  (6.95,1.00) -- (5.30,1.00) -- (5.30,1.35) -- (5.90,1.95) -- (5.30,2.55) -- (5.30,2.90) -- (6.95,2.90) -- cycle;
\node[num, text=solideG] at (6.50,1.95) {7};

%% ---------------- 3 : la vis (filet à gauche puis filet à droite) --------------------------
\draw[piece, draw=solideB, fill=white] (1.75,0.34) rectangle (9.15,0.66);
\foreach \x in {1.90,2.14,...,4.60} { \draw[line width=0.5pt, solideB] (\x,0.66) -- (\x+0.16,0.34); }
\foreach \x in {4.80,5.04,...,9.00} { \draw[line width=0.5pt, solideB] (\x,0.34) -- (\x+0.16,0.66); }
\draw[piece, draw=solideB] (4.70,0.30) -- (4.70,0.70);
\node[num, text=solideB] at (4.98,0.50) {3};

% le tube, posé entre les deux V
\draw[piece, draw=solideD, fill=solideD!10] (4.70,1.95) circle (0.85);
\draw[piece, draw=solideD, line width=0.6pt] (4.70,1.95) circle (0.66);
\node[num, text=solideD] at (4.70,1.95) {T};

%% ---------------- les deux sens de filetage (donnée de l'énoncé) ---------------------------
\draw[renvoi] (1.55,1.78) -- (2.05,0.70);
\node[etiq, anchor=south, align=center, text=solideB] at (1.50,1.80) {pas à\\gauche};
\draw[renvoi] (8.05,1.78) -- (7.55,0.70);
\node[etiq, anchor=south, align=center, text=solideB] at (8.10,1.80) {pas à\\droite};

%% ---------------- les mors se rapprochent (donnée de l'énoncé) -----------------------------
\draw[fleche, solideC] (2.30,3.15) -- (2.95,3.15);
\draw[fleche, solideG] (7.10,3.15) -- (6.45,3.15);
\node[etiq, text=black!75] at (4.70,3.15) {les mors se rapprochent};

%% ---------------- repère lié au bâti, $\vec x$ suivant l'axe de la vis ---------------------
\draw[-{Stealth[length=5pt]}, line width=0.6pt] (0.62,-1.30) -- ++(0.70,0)
  node[anchor=west, inner sep=2pt, font=\small] {$\vec x$};
\draw[-{Stealth[length=5pt]}, line width=0.6pt] (0.62,-1.30) -- ++(0,0.42)
  node[anchor=west, inner sep=2pt, font=\small] {$\vec y$};
\fill (0.62,-1.30) circle (1.3pt);
\node[etiq, anchor=east, inner sep=3pt] at (0.62,-1.30) {$O$};

%% ---------------- la réponse : les classes d'équivalence -----------------------------------
\rappel{\node[draw=black!45, fill=black!3, rounded corners=3pt, inner sep=5pt, anchor=north,
  font=\small] at (5.55,-0.86)
  {classes : $E_0=\{1,2\}$\; $E_1=\{3\}$\; $E_2=\{4,5\}$\; $E_3=\{6,7\}$};}
\end{tikzpicture}
\end{document}
